\chapter{可计算性}
\label{chap:computability}

上一章定义了图灵机。本章研究\emph{可计算性}（computability）：首先说明一台
图灵机可以模拟任意其他图灵机，然后研究哪些函数无法由任何图灵机计算。

\section{通用图灵机}

为了把机器本身作为输入，需要先约定图灵机的编码。对
$M=(K,\Gamma,s,\delta)$，分别把 $K$ 与 $\Gamma$ 中的元素编号，再把每条转移

\[
  \delta(q,a)=(q',b,D)
\]

写成五元组 $(q,a,q',b,D)$。利用自定界编码依次记录状态数、字母表大小、初始
状态和全部转移，就得到 $M$ 的编码 $\langle M\rangle\in\bitsstar$。这种编码
具有以下性质：是否合法可以判定，并且可以从合法编码中恢复 $M$ 的全部数据。

下文用 $\langle u,v\rangle$ 表示两个字符串的固定自定界编码，并把
$U(\langle m,x\rangle)$ 简写为 $U(m,x)$。

\begin{theorem}[通用图灵机]
  \label{thm:universal-turing-machine}
  存在一台图灵机 $U$，使得对任意 $m,x\in\bitsstar$，

  \[
    U(m,x)=
    \begin{cases}
      M(x),&m=\langle M\rangle\text{ 是某台图灵机的合法编码},\\
      0,&m\text{ 不是合法编码}.
    \end{cases}
  \]

  第一种情形中，若 $M$ 在 $x$ 上不停机，则 $U$ 也不停机。
\end{theorem}

\begin{proof}
  $U$ 先检查并解析 $m$。若编码不合法，就输出 $0$；否则从中恢复
  $M=(K,\Gamma,s,\delta)$。随后 $U$ 在自己的纸带上记录 $M$ 的当前状态、读写
  头位置和纸带内容，并根据编码中的转移表逐步更新这三个部分。每一轮模拟 $M$
  的一步，所以两台机器同时停机且输出相同。
\end{proof}

这样的 $U$ 称为\emph{通用图灵机}（universal Turing machine）。它表明我们
不需要为每个算法制造一台新的物理机器：只需把程序编码后交给同一台机器执行。

\section{不可计算函数}

\begin{theorem}[不可计算函数的存在性]
  \label{thm:uncomputable-function-exists}
  存在函数 $F:\bitsstar\to\bits$ 不是可计算函数。
\end{theorem}

\begin{proof}[证明一：计数]
  每台图灵机都有一个有限二进制编码，所以图灵机的集合至多可数，因而它们至多
  能计算可数多个函数。另一方面，$\bitsstar$ 是可数无限集，因此从
  $\bitsstar$ 到 $\bits$ 的函数共有

  \[
    2^{\aleph_0}
  \]

  个，是不可数的。故其中必有函数不能由图灵机计算。
\end{proof}

下面的对角化（diagonalization）证明给出一个具体的不可计算函数。若 $m$ 是
合法图灵机编码，记其对应的机器为 $M_m$。定义

\begin{equation}
  \label{eq:diagonal-uncomputable-function}
  F^*(m)=
  \begin{cases}
    0,&m\text{ 合法且 }M_m(m)=1,\\
    1,&\text{其他情形}.
  \end{cases}
\end{equation}

“其他情形”包括 $m$ 非法、$M_m(m)$ 不停机，以及它停机但输出不是 $1$。

\begin{proof}[证明二：对角化]
  假设某台图灵机 $M_e$ 计算 $F^*$，其中 $e=\langle M_e\rangle$。把 $e$ 输入
  它自身。若 $M_e(e)=1$，则由 \cref{eq:diagonal-uncomputable-function} 得
  $F^*(e)=0$；若 $M_e(e)=0$，则同一式给出 $F^*(e)=1$。两种情形都与
  $M_e(e)=F^*(e)$ 矛盾，故 $F^*$ 不可计算。
\end{proof}

\section{停机问题}

\begin{definition}[停机问题]
  \label{def:halting-problem}
  \emph{停机问题}（halting problem）是函数

  \[
    \operatorname{HALT}(m,x)=
    \begin{cases}
      1,&m=\langle M\rangle\text{ 且 }M\text{ 在输入 }x\text{ 上停机},\\
      0,&\text{其他情形}.
    \end{cases}
  \]
\end{definition}

\begin{theorem}[停机问题不可计算]
  \label{thm:halting-problem-uncomputable}
  函数 $\operatorname{HALT}$ 不可计算。
\end{theorem}

\begin{proof}
  假设图灵机 $H$ 计算 $\operatorname{HALT}$。由它构造一台机器 $D$：输入
  $m$ 后，先计算 $H(m,m)$。若结果为 $0$，则输出 $1$；若结果为 $1$，则按照
  \cref{thm:universal-turing-machine} 的方法逐步模拟 $M_m(m)$ 直至其停机。
  若模拟得到的输出恰为 $1$，就输出 $0$；否则输出 $1$。

  \begin{pseudocode}{Diagonal machine $D$ (input $m$)}
    \item Compute $h\gets H(m,m)$.
    \item If $h=0$, output $1$ and halt.
    \item Otherwise, simulate $M_m(m)$ until it halts.
    \item If the simulation outputs $1$, output $0$; otherwise, output $1$.
  \end{pseudocode}

  当 $H(m,m)=1$ 时，被模拟的机器必定停机，所以 $D$ 在每个输入上都停机。
  对照 \cref{eq:diagonal-uncomputable-function} 可知 $D$ 恰好计算 $F^*$，这与
  \cref{thm:uncomputable-function-exists} 矛盾。
\end{proof}

只询问机器在某个固定输入上是否停机，也不会使问题变得可计算。定义

\[
  \operatorname{HALT}_0(m)=
  \begin{cases}
    1,&m=\langle M\rangle\text{ 且 }M\text{ 在输入 }0\text{ 上停机},\\
    0,&\text{其他情形}.
  \end{cases}
\]

\begin{theorem}[固定输入上的停机问题]
  \label{thm:halt-on-zero-uncomputable}
  函数 $\operatorname{HALT}_0$ 不可计算。
\end{theorem}

\begin{proof}
  给定 $m,x$，可以有效构造图灵机 $N_{m,x}$：它忽略自己的输入，若 $m$ 合法
  就模拟 $M_m(x)$，否则永远循环。于是

  \begin{pseudocode}{Machine $N_{m,x}$ (input $y$)}
    \item Ignore $y$.
    \item If $m$ is not a valid machine encoding, loop forever.
    \item Otherwise, simulate $M_m(x)$.
  \end{pseudocode}

  \[
    \operatorname{HALT}(m,x)
      =\operatorname{HALT}_0(\langle N_{m,x}\rangle).
  \]

  若某台机器 $Z$ 计算 $\operatorname{HALT}_0$，等式右侧的算法可写成：

  \begin{pseudocode}{Halting decider $H$ (input $\langle m,x\rangle$)}
    \item Construct $N_{m,x}$ and its encoding $n\gets\langle N_{m,x}\rangle$.
    \item Run $Z(n)$ and output its result.
  \end{pseudocode}

  这便给出了停机问题的算法，与
  \cref{thm:halting-problem-uncomputable} 矛盾。
\end{proof}

不可计算性不仅涉及机器是否停机，也涉及机器计算了什么函数。

\begin{example}[识别常零函数]
  \label{ex:zero-function-uncomputable}
  令 $Z(x)=0$，并定义

  \[
    \operatorname{ZERO}(m)=
    \begin{cases}
      1,&m=\langle M\rangle\text{ 且 }M\text{ 计算函数 }Z,\\
      0,&\text{其他情形}.
    \end{cases}
  \]

  函数 $\operatorname{ZERO}$ 不可计算。
\end{example}

\begin{proof}
  假设某台机器 $E$ 计算 $\operatorname{ZERO}$。给定 $m,x$，构造机器
  $N_{m,x}$：在任意
  输入 $y$ 上，它先模拟 $M_m(x)$；若模拟停机，就输出 $0$。若 $m$ 非法，则令
  $N_{m,x}$ 永远循环。

  \begin{pseudocode}{Machine $N_{m,x}$ (input $y$)}
    \item Ignore $y$.
    \item If $m$ is not a valid machine encoding, loop forever.
    \item Simulate $M_m(x)$ until it halts.
    \item Output $0$ and halt.
  \end{pseudocode}

  因此

  \[
    \operatorname{HALT}(m,x)
      =\operatorname{ZERO}(\langle N_{m,x}\rangle),
  \]

  且停机问题的判定器可写成：

  \begin{pseudocode}{Halting decider $H$ (input $\langle m,x\rangle$)}
    \item Construct $N_{m,x}$ and its encoding $n\gets\langle N_{m,x}\rangle$.
    \item Run $E(n)$ and output its result.
  \end{pseudocode}

  这将使停机问题可计算，产生矛盾。
\end{proof}

\section{归约}

\begin{definition}[归约]
  \label{def:computable-reduction}
  设 $F,G:\bitsstar\to\bits$。从 $F$ 到 $G$ 的一个\emph{归约}
  （reduction）是一个可计算函数 $R:\bitsstar\to\bitsstar$，满足

  \[
    F(x)=G(R(x))
    \qquad(x\in\bitsstar).
  \]

  若存在这样的 $R$，记为 $F\leq_{\mathrm m}G$，并称 $F$ 可多对一归约到 $G$
  （many-one reducible）。
\end{definition}

\begin{corollary}
  \label{cor:reduction-preserves-computability}
  若 $F\leq_{\mathrm m}G$ 且 $G$ 可计算，则 $F$ 可计算。
\end{corollary}

\begin{proof}
  先计算 $R(x)$，再计算 $G(R(x))$，所得结果就是 $F(x)$。
\end{proof}

\begin{corollary}
  \label{cor:reduction-transfers-uncomputability}
  若 $F\leq_{\mathrm m}G$ 且 $F$ 不可计算，则 $G$ 不可计算。
\end{corollary}

\begin{proof}
  这是 \cref{cor:reduction-preserves-computability} 的逆否命题。
\end{proof}

前面两个证明已经隐含使用了归约：它们分别给出

\[
  \operatorname{HALT}\leq_{\mathrm m}\operatorname{HALT}_0,
  \qquad
  \operatorname{HALT}\leq_{\mathrm m}\operatorname{ZERO}.
\]

\section{Rice 定理}

机器的\emph{语义性质}（semantic property）只取决于它所计算的函数，而不
取决于状态名称、转移表写法等实现细节。

\begin{definition}[语义且非平凡的性质]
  \label{def:semantic-nontrivial-property}
  图灵机的性质 $P$ 称为\emph{语义的}，如果任意两台计算同一个偏函数
  （partial function）的机器 $M,M'$ 都满足

  \[
    P(M)=P(M').
  \]

  若存在图灵机 $M_0,M_1$ 使得 $P(M_0)=0$ 且 $P(M_1)=1$，则称 $P$ 是
  \emph{非平凡的}（non-trivial）。
\end{definition}

对这样的性质，考虑其判定函数

\[
  \chi_P(m)=
  \begin{cases}
    P(M),&m=\langle M\rangle,\\
    0,&m\text{ 不是合法编码}.
  \end{cases}
\]

\begin{theorem}[Rice 定理]
  \label{thm:rice}
  若 $P$ 是图灵机的语义且非平凡的性质，则 $\chi_P$ 不可计算。
\end{theorem}

\begin{proof}
  令 $M_\infty$ 为在每个输入上都永远循环的机器。必要时把 $P$ 换成其补性质
  $1-P$；判定二者是等价的。因此不妨设 $P(M_\infty)=0$。由非平凡性可取一台
  机器 $M_{\mathrm{yes}}$，使得 $P(M_{\mathrm{yes}})=1$。

  给定 $m,x$，有效构造机器 $N_{m,x}$。它在输入 $y$ 上执行：

  \begin{enumerate}[label=\arabic*.]
    \item 若 $m$ 非法则永远循环；否则模拟 $M_m(x)$；
    \item 上一步停机后，模拟 $M_{\mathrm{yes}}(y)$ 并原样输出其结果。
  \end{enumerate}

  若 $\operatorname{HALT}(m,x)=0$，则 $N_{m,x}$ 与 $M_\infty$ 计算同一个处处
  无定义的偏函数，所以 $P(N_{m,x})=0$。若
  $\operatorname{HALT}(m,x)=1$，则 $N_{m,x}$ 与 $M_{\mathrm{yes}}$ 计算同一个
  偏函数，所以 $P(N_{m,x})=1$。因此

  \[
    \operatorname{HALT}(m,x)
      =\chi_P(\langle N_{m,x}\rangle),
  \]

  即 $\operatorname{HALT}\leq_{\mathrm m}\chi_P$。由
  \cref{thm:halting-problem-uncomputable} 以及
  \cref{cor:reduction-transfers-uncomputability} 可知 $\chi_P$ 不可计算。
\end{proof}

\lecture
\label{lec:recursion-and-incompleteness}

\section{自引用与递归定理}

程序可以把程序编码作为普通字符串处理。固定上一节的图灵机编码后，存在以下
两个可计算的代码生成函数：

\begin{itemize}
  \item $q(w)$ 返回机器 $A_w$ 的编码；$A_w$ 在输入 $x$ 上输出二元组
    $\langle w,x\rangle$；
  \item $c(a,b)$ 返回顺次运行编码为 $a,b$ 的两台机器所得机器的编码。
\end{itemize}

它们都是可计算的，因为只需把 $w$ 或两张转移表填入一个固定的程序模板。

\begin{pseudocode}{Hardwiring machine $A_w$ (input $x$)}
  \item Preserve $x$ and write the fixed string $w$ on a work tape.
  \item Output $\langle w,x\rangle$ and halt.
\end{pseudocode}

\begin{example}[自打印图灵机]
  \label{ex:self-printing-tm}
  存在一台图灵机 $R$，它在任意输入上都输出自己的编码
  $\langle R\rangle$。
\end{example}

\begin{proof}
  构造机器 $B$：它在输入 $\langle w,x\rangle$ 上输出 $c(q(w),w)$。令

  \begin{pseudocode}{Machine $B$ (input $\langle w,x\rangle$)}
    \item Compute $a\gets q(w)$, the encoding of $A_w$.
    \item Compute $r_w\gets c(a,w)$.
    \item Output $r_w$ and halt.
  \end{pseudocode}

  \[
    b=\langle B\rangle,
    \qquad
    r=c(q(b),b),
  \]

  并令 $R$ 为编码 $r$ 所表示的机器。$R$ 先运行 $A_b$，把输入 $x$ 变为
  $\langle b,x\rangle$，再运行 $B$，所以

  \begin{pseudocode}{Self-printing machine $R$ (input $x$)}
    \item Run $A_b(x)$ to obtain $\langle b,x\rangle$.
    \item Run $B$ on the result and return its output.
  \end{pseudocode}

  \[
    R(x)=c(q(b),b)=r=\langle R\rangle.
  \]
\end{proof}

同一个技巧允许程序在运行时取得自己的编码。

\begin{theorem}[递归定理]
  \label{thm:recursion-theorem}
  设图灵机 $T$ 接受两个输入。存在图灵机 $R$，使得对每个
  $x\in\bitsstar$，

  \[
    R(x)=T(\langle R\rangle,x).
  \]

  等式按偏函数理解：一侧不停机当且仅当另一侧不停机。
\end{theorem}

\begin{proof}
  由 $T$ 构造机器 $B_T$。它在输入 $\langle w,x\rangle$ 上先计算

  \[
    r_w=c(q(w),w),
  \]

  再运行 $T(r_w,x)$。令 $b=\langle B_T\rangle$，并令 $R$ 是编码
  $r=c(q(b),b)$ 所表示的机器。对任意 $x$，第一阶段 $A_b$ 产生
  $\langle b,x\rangle$，随后 $B_T$ 计算的 $r_b$ 正好等于
  $r=\langle R\rangle$。因此

  \begin{pseudocode}{Machine $B_T$ in the recursion theorem
      (input $\langle w,x\rangle$)}
    \item Compute $a\gets q(w)$.
    \item Compute $r_w\gets c(a,w)$.
    \item Simulate $T(r_w,x)$ and return its output.
  \end{pseudocode}

  \begin{pseudocode}{Fixed-point machine $R$ (input $x$)}
    \item Run $A_b(x)$ to obtain $\langle b,x\rangle$, where
      $b=\langle B_T\rangle$.
    \item Run $B_T$ on the result and return its output.
  \end{pseudocode}

  \[
    R(x)=B_T(b,x)=T(\langle R\rangle,x).
  \]
\end{proof}

\begin{example}[用递归定理重证停机问题]
  \label{ex:halting-via-recursion-theorem}
  假设存在计算 $\operatorname{HALT}$ 的机器 $H$。构造机器 $T$，使得

  \[
    T(m,x)=
    \begin{cases}
      \text{永远循环},&H(m,x)=1,\\
      0,&H(m,x)=0.
    \end{cases}
  \]

  等价地，由递归定理得到的自引用机器可写成：

  \begin{pseudocode}{Contradiction machine $R$ (input $x$)}
    \item Obtain its own encoding $m\gets\langle R\rangle$.
    \item Compute $h\gets H(m,x)$.
    \item If $h=1$, loop forever; otherwise, output $0$ and halt.
  \end{pseudocode}

  由 \cref{thm:recursion-theorem}，存在 $R$ 满足
  $R(x)=T(\langle R\rangle,x)$。若 $H(\langle R\rangle,x)=1$，则 $R(x)$
  不停机；若其值为 $0$，则 $R(x)$ 停机。两种情形都与 $H$ 的输出矛盾，故
  $\operatorname{HALT}$ 不可计算。
\end{example}

用老师的话来说，有了递归定理，就可以通过这种“作弊”的方式来证明。

\subsection{最小机器与不动点}

记 $M\simeq N$ 表示两台机器计算同一个偏函数。称机器 $M$ 是\emph{最小的}
（minimal），如果

\[
  M'\simeq M
  \quad\Longrightarrow\quad
  \card{\langle M'\rangle}\geq\card{\langle M\rangle}.
\]

定义

\[
  \operatorname{MIN}(m)=
  \begin{cases}
    1,&m=\langle M\rangle\text{ 且 }M\text{ 是最小机器},\\
    0,&\text{其他情形}.
  \end{cases}
\]

\begin{example}[识别最小机器]
  \label{ex:minimal-tm-uncomputable}
  函数 $\operatorname{MIN}$ 不可计算。
\end{example}

\begin{proof}
  假设 $\operatorname{MIN}$ 可计算。最小机器的编码长度没有上界：函数
  $x\mapsto 0^n$ 随 $n$ 改变而彼此不同，而长度不超过任意给定常数的编码只有
  有限多个。

  构造机器 $T$。在输入 $\langle m,x\rangle$ 上，它按长度枚举字符串 $s$，直至
  找到

  \[
    \operatorname{MIN}(s)=1
    \quad\text{且}\quad
    \card{s}>\card{m},
  \]

  然后在 $x$ 上模拟编码为 $s$ 的机器。由无界性，这个搜索一定结束。根据
  \cref{thm:recursion-theorem}，存在 $R$ 使得
  $R(x)=T(\langle R\rangle,x)$。设搜索得到 $s$，则

  \begin{pseudocode}{Contradiction machine $R$ for recognizing minimal
      machines (input $x$)}
    \item Obtain its own encoding $m\gets\langle R\rangle$.
    \item Enumerate $s\in\bitsstar$ in order of length until
      $\operatorname{MIN}(s)=1$ and $\card{s}>\card{m}$.
    \item Simulate $M_s$ on $x$ and return its output.
  \end{pseudocode}

  \[
    R\simeq M_s,
    \qquad
    \card{\langle R\rangle}<\card{s}.
  \]

  这与 $M_s$ 是最小机器矛盾。
\end{proof}

递归定理也常写成下面的不动点（fixed point）形式。

\begin{corollary}[不动点定理]
  \label{cor:fixed-point-theorem}
  设 $f:\bitsstar\to\bitsstar$ 是可计算函数，并且把合法图灵机编码映射为合法
  图灵机编码。则存在图灵机 $F$，使得

  \[
    F\simeq M_{f(\langle F\rangle)}.
  \]
\end{corollary}

\begin{proof}
  令 $T(m,x)$ 模拟编码为 $f(m)$ 的机器在 $x$ 上的计算。对 $T$ 应用
  \cref{thm:recursion-theorem} 即得所需的 $F$。

  \begin{pseudocode}{Fixed-point machine $F$ (input $x$)}
    \item Obtain its own encoding $e\gets\langle F\rangle$.
    \item Compute $m\gets f(e)$.
    \item Simulate $M_m$ on $x$ and return its output.
  \end{pseudocode}
\end{proof}

\section{有效证明系统与不完备性}

\begin{definition}[有效证明系统]
  \label{def:effective-proof-system}
  设 $\mathcal T\subseteq\bitsstar$ 是一组真命题的编码。它的一个\emph{有效
  证明系统}（effective proof system）是图灵机

  \[
    V:\bitsstar\times\bitsstar\to\bits,
  \]

  其中 $V(x,p)=1$ 表示 $p$ 是命题 $x$ 的证明，并满足：

  \begin{enumerate}[label=\arabic*.]
    \item \textbf{有效性（effectiveness）：}对每个 $x,p$，$V(x,p)$ 都停机；
    \item \textbf{可靠性（soundness）：}若 $x\notin\mathcal T$，则对每个 $p$
      都有 $V(x,p)=0$。
  \end{enumerate}

  如果对每个 $x\in\mathcal T$ 都存在 $p$ 使 $V(x,p)=1$，则称 $V$ 是
  \emph{完备的}（complete）。
\end{definition}

\begin{theorem}[抽象不完备性]
  \label{thm:abstract-incompleteness}
  存在一组真命题 $\mathcal T$，它没有可靠且完备的有效证明系统。
\end{theorem}

\begin{proof}
  对每台图灵机 $M$，考虑两个命题

  \[
    h_M=\text{“$M$ 在输入 $0$ 上停机”},
    \qquad
    n_M=\text{“$M$ 在输入 $0$ 上不停机”},
  \]

  并令 $\mathcal T$ 包含其中所有真命题。假设 $V$ 是可靠且完备的有效证明系统。
  给定字符串 $m$，若它不是合法机器编码就输出 $0$；否则写成
  $m=\langle M\rangle$，并同时进行以下两项计算：

  \begin{enumerate}[label=\arabic*.]
    \item 逐步运行 $M(0)$；若它停机，则输出 $1$；
    \item 枚举所有 $p\in\bitsstar$ 并计算 $V(n_M,p)$；若某次得到 $1$，则输出
      $0$。
  \end{enumerate}

  \begin{pseudocode}{Halting decider $D$ from a complete proof system
      (input $m$)}
    \item If $m$ is not a valid Turing-machine encoding, output $0$;
      otherwise, decode it as $M$.
    \item Dovetail two branches: simulate $M(0)$ step by step, and enumerate
      $p\in\bitsstar$ while computing $V(n_M,p)$.
    \item If the first branch finds that $M(0)$ halts, output $1$; if the
      second finds $V(n_M,p)=1$, output $0$.
  \end{pseudocode}

  若 $M(0)$ 停机，第一项计算会结束；若它不停机，则 $n_M\in\mathcal T$，完备性
  保证第二项最终找到证明。可靠性保证两项不会给出冲突的答案。因此这个过程计算
  $\operatorname{HALT}_0$，与 \cref{thm:halt-on-zero-uncomputable} 矛盾。
\end{proof}

这个结论给出了哥德尔不完备现象的计算核心。为了把它落实到算术语句，下面固定
一种具体的形式语言。

\subsection{量化整数语句}

\begin{definition}[量化整数语句]
  \label{def:quantified-integer-statement}
  \emph{量化整数语句}（quantified integer statement, QIS）是由下列成分组成的
  闭公式：

  \begin{itemize}
    \item 非负整数常数、变量以及运算 $+$、$\times$ 和关系 $=$；
    \item 逻辑联结词 $\land,\lor,\lnot$；
    \item 量词 $\forall,\exists$。
  \end{itemize}

  本节约定量词的论域为 $\Nat$。闭公式没有自由变量，因而在标准自然数结构中
  具有确定的真值。
\end{definition}

对非法编码约定结果为 $0$，并定义

\[
  \operatorname{EVAL}(\varphi)=
  \begin{cases}
    1,&\varphi\text{ 是真的 QIS},\\
    0,&\varphi\text{ 是假的 QIS 或不是合法 QIS}.
  \end{cases}
\]

\begin{lemma}[计算历史的算术化]
  \label{lem:arithmetization-of-computation}
  存在一个可计算变换，把任意图灵机编码 $\langle M\rangle$ 映射为 QIS
  $\varphi_M$，使得

  \[
    \varphi_M\text{ 为真}
    \quad\Longleftrightarrow\quad
    M\text{ 在输入 }0\text{ 上停机}.
  \]
\end{lemma}

\begin{proof}[证明思路]
  用自然数编码状态、纸带符号和有限格局，再用一个自然数编码有限格局序列。
  “后一格局由前一格局经一条转移得到”是有限的算术条件；序列取位、长度和相邻
  关系也可用加法、乘法和量词表达。因此可以有效写出一个公式，断言“存在一个从
  $M$ 的初始格局开始、以停机格局结束的有限合法计算历史”。
\end{proof}

\begin{theorem}[QIS 真值不可判定]
  \label{thm:qis-evaluation-uncomputable}
  函数 $\operatorname{EVAL}$ 不可计算。
\end{theorem}

\begin{proof}
  若 $\operatorname{EVAL}$ 可计算，则由
  \cref{lem:arithmetization-of-computation} 可得

  \begin{pseudocode}{Halting decider $D_0$ using $\operatorname{EVAL}$
      (input $m$)}
    \item If $m$ is not a valid Turing-machine encoding, output $0$.
    \item Use the arithmetization lemma to compute $\varphi_{M_m}$.
    \item Output $\operatorname{EVAL}(\varphi_{M_m})$.
  \end{pseudocode}

  \[
    \operatorname{HALT}_0(\langle M\rangle)
      =\operatorname{EVAL}(\varphi_M),
  \]

  这与 \cref{thm:halt-on-zero-uncomputable} 矛盾。
\end{proof}

\begin{theorem}[量化整数真理的不完备性]
  \label{thm:qis-incompleteness}
  所有真 QIS 构成的集合没有可靠且完备的有效证明系统。
\end{theorem}

\begin{proof}
  假设存在这样的验证器 $V$。给定合法 QIS $\varphi$，同时枚举
  $\varphi$ 与 $\lnot\varphi$ 的所有候选证明，并用 $V$ 验证。二者恰有一个为真；
  完备性保证其证明最终被找到，可靠性保证假命题不会被证明。因此我们可以判定
  $\varphi$ 的真值，这与 \cref{thm:qis-evaluation-uncomputable} 矛盾。

  \begin{pseudocode}{QIS decider $D_{\mathrm{QIS}}$ from a complete proof
      system (input $\varphi$)}
    \item If $\varphi$ is not a valid QIS, output $0$.
    \item Dovetail candidate proofs $p$ for $\varphi$ and $\lnot\varphi$,
      evaluating $V(\varphi,p)$ and $V(\lnot\varphi,p)$, respectively.
    \item If a proof of $\varphi$ is found, output $1$; if a proof of
      $\lnot\varphi$ is found, output $0$.
  \end{pseudocode}
\end{proof}
